% Generated draft: requires mathematical and source review.
\documentclass[11pt]{article}
\usepackage[margin=25mm]{geometry}
\usepackage{fontspec}
\setmainfont{Latin Modern Roman}
\usepackage{amsmath,amssymb,mathtools,unicode-math}
\setmathfont{Latin Modern Math}
\usepackage{enumitem,needspace,xurl,hyperref}
\setlength{\parindent}{0pt}
\setlength{\parskip}{5pt}
\setlength{\emergencystretch}{3em}
\newcommand{\N}{\mathbb{N}}
\newcommand{\Z}{\mathbb{Z}}
\newcommand{\Q}{\mathbb{Q}}
\newcommand{\R}{\mathbb{R}}
\newcommand{\C}{\mathbb{C}}
\newcommand{\F}{\mathbb{F}}
\newcommand{\A}{\mathbb{A}}
\newcommand{\PP}{\mathbb{P}}
\newcommand{\E}{\mathbb{E}}
\newcommand{\eps}{\varepsilon}
\newcommand{\dd}{\,\mathrm{d}}
\begin{document}
\section*{UnsolvedMath ID: 30004158}
\textbf{Source problem code:} catalogue entry
\subsection{Definitions and mathematical statement}
Let \((M,D)\) be a pair described in the supplied source as positive and maximally degenerate, and put \(X=M\setminus D\). Let \(A\) denote the ring generated by the Gross--Siebert theta functions associated with this pair. Write \(\star_A\) for its multiplication and \(\star_{SH}\) for the product on the symplectic cohomology \(SH^*(X)\), with the geometric and coefficient choices required for these constructions understood to be fixed consistently.

The proposed assertion is that, for every pair in the intended class, there exists a ring isomorphism
\[
\Phi:A\xrightarrow{\sim}SH^*(X).
\]
Thus \(\Phi\) should be an additive bijection satisfying
\[
\Phi(a\star_A b)=\Phi(a)\star_{SH}\Phi(b)
\qquad\text{for all }a,b\in A.
\]
If the intended ring structures and isomorphism are unital, one additionally requires \(\Phi(1_A)=1_{SH}\).

\textit{Unresolved qualifications.} The supplied record does not define positivity or maximal degeneracy, specify the dimension or geometric category of the pair, or give the hypotheses and auxiliary structures needed to define these particular theta functions and symplectic cohomology. It also leaves unspecified the coefficient rings, any completion or specialization, and whether the comparison preserves a grading or a coefficient algebra structure. Accordingly, the display records the intended assertion schematically; these missing specifications prevent a fully determined mathematical formulation. The notation \(SH^*(X)\) is retained from the source and does not resolve whether all degrees or only a specified degree are intended.

The supplied edited statement and literature assessment describe an established additive identification. The original statement instead refers to cases where an additive identification is known. Neither gives its precise scope or specifies whether \(\Phi\) must extend a particular additive identification. Compatibility of a prescribed additive identification with multiplication is a more specific target than existence of some ring isomorphism; the record does not settle which is intended.
\subsection{Short English statement}
For a positive maximally degenerate pair \((M,D)\), the question asks whether its Gross--Siebert theta-function ring is isomorphic to \(SH^*(M\setminus D)\) as a ring. The intended geometric hypotheses, coefficients, grading conventions, and relationship to the reported additive identification require clarification.
\subsection{Sources}
\begin{itemize}
\item Local UnsolvedMath registry: permanent ID \(30004158\), title \emph{Theta Function Rings and Symplectic Cohomology}, source problem code \texttt{catalogue entry}, in the Oberwolfach Reports Open Problems collection. No individual proposer is supplied; the proposed year is recorded as 2019.
\item Supplied report citation: \emph{Logarithmic Enumerative Geometry and Mirror Symmetry} (2020), Oberwolfach Reports, pp. 1639--1695, DOI \texttt{10.4171/owr/2019/27}. Listed also as the OWR primary report: \url{https://doi.org/10.4171/owr/2019/27}.
\item Supplied literature reference, labeled ``Pascaleff on symplectic cohomology of log Calabi--Yau surfaces'': \url{https://doi.org/10.2140/gt.2019.23.2701}.
\item Literature assessment dated exactly \texttt{2026-08-22}: the supplied curation labels the problem partially solved and reports an established additive identification with theta-function basis elements, product computations in important special cases, and a ring-isomorphism conjecture remaining in the stated generality. This is supplied dated curation derived from the OWR record, not an independently verified literature assessment; its status and cited claims require independent verification.
\item The registry labels the statement \texttt{corrected\_verified} and describes repairs to notation, typography, and sentence structure while preserving mathematical content. This supplied editorial label does not establish mathematical correctness or resolve the missing hypotheses above.
\end{itemize}
\end{document}
